\section*{Capítulo \ref{ch:b2:cell-differentiation} --- La diferenciación celular: la célula muscular esquelética}

\begin{solution}{exo:b2:cell-differentiation:1}
\omterm{def:b2:cell-differentiation:differentiation}{Diferenciación}: la expresión estable de los genes y las funciones de un
tipo celular. \omterm{def:b2:cell-differentiation:differentiation}{Determinación}: el compromiso con un destino antes de que se
exprese. \omterm{def:b2:cell-differentiation:differentiation}{Potencia}: el abanico de destinos todavía abiertos. Cigoto:
totipotente; célula de la \omterm{def:b2:mammal-reproduction:implantation}{masa celular interna}: pluripotente; célula
satélite: unipotente (solo músculo) --- aun así, una \omterm{def:b2:cell-differentiation:differentiation}{célula madre}; \omterm{def:b2:cell-differentiation:myogenesis}{fibra
muscular}: diferenciada, posmitótica, sin \omterm{def:b2:cell-differentiation:differentiation}{potencia}.
\end{solution}

\begin{solution}{exo:b2:cell-differentiation:2}
Célula del dermomiotomo (Pax3) $\to$ \omterm{def:b2:cell-differentiation:myogenesis}{mioblasto} inducido (\omterm{def:b2:cell-differentiation:myogenesis}{MyoD}, Myf5),
que se divide mientras dura el FGF $\to$ salida del ciclo, miogenina $\to$
alineamiento y fusión en un \omterm{def:b2:cell-differentiation:myogenesis}{miotubo} $\to$ genes musculares activos,
\omterm{def:b2:cell-differentiation:fibre}{miofibrillas} ensambladas (MRF4) $\to$ fibra madura con núcleos
periféricos, más células satélite en reposo (Pax7).
\end{solution}

\begin{solution}{exo:b2:cell-differentiation:3}
Que el núcleo de una célula diferenciada conserva todos los genes
necesarios para formar un animal completo, y que el citoplasma del huevo
puede reiniciar su programa: la \omterm{def:b2:cell-differentiation:differentiation}{diferenciación} es reversible y no es una
pérdida de ADN. No mostró que todo núcleo diferenciado pueda
reiniciarse (la mayoría de los trasplantes fracasaron), ni cómo funciona
el reinicio, ni que la propia célula diferenciada pudiera cambiar de
destino allí donde estaba.
\end{solution}

\begin{solution}{exo:b2:cell-differentiation:4}
Los discos Z que delimitan el \omterm{def:b2:cell-differentiation:fibre}{sarcómero}; los filamentos finos de actina
anclados en ellos; los filamentos gruesos de miosina centrados en la línea
M; la banda A (filamentos gruesos) y la banda I (solo filamentos finos).
Al acortarse, los filamentos finos se deslizan hacia la línea M y los
discos Z se acercan entre sí: la banda I y la zona H se estrechan, la
banda A --- la longitud de los filamentos gruesos --- no cambia, y ningún
filamento se acorta.
\end{solution}

\begin{solution}{exo:b2:cell-differentiation:5}
$0.2/2.5\times 10^{-6} = \num{80000}$ \omterm{def:b2:cell-differentiation:fibre}{sarcómeros}; $0.2/25\times 10^{-6}
= \num{8000}$ núcleos; un \omterm{def:b2:cell-differentiation:myogenesis}{mioblasto} por núcleo, de modo que se fusionaron
unos \num{8000} \omterm{def:b2:cell-differentiation:myogenesis}{mioblastos} (las células satélite añaden más después).
\end{solution}

\begin{solution}{exo:b2:cell-differentiation:6}
El FGF induce Id, que se une a la proteína compañera que \omterm{def:b2:cell-differentiation:myogenesis}{MyoD} necesita
para formar un dímero capaz de unirse al ADN; \omterm{def:b2:cell-differentiation:myogenesis}{MyoD} está presente pero
inactivo, y las células siguen dividiéndose. Cuando se retira el FGF, Id
se degrada, \omterm{def:b2:cell-differentiation:myogenesis}{MyoD} dimeriza, se activa la miogenina y las células salen del
ciclo y se fusionan. Los \omterm{def:b2:cell-differentiation:myogenesis}{mioblastos} que fabrican Id de forma constitutiva
nunca se fusionan: se dividen indefinidamente y no forman músculo.
\end{solution}

\begin{solution}{exo:b2:cell-differentiation:7}
Condición $0.4m = m^{2}/(1 + m^{2}) + 0.02$. Con $m = 2.1$: miembro
izquierdo $0.84$, derecho $4.41/5.41 + 0.02 = 0.835$: un estado
estacionario. Estado bajo: para $m$ pequeña, $0.4m \approx m^{2} + 0.02$,
luego $m \approx
0.055$ (izquierdo $0.022$, derecho $0.023$). Elevada a 1, por encima del
umbral situado cerca de $0.41$, la célula sube hasta $2.1$ y se queda ahí:
está comprometida.
\end{solution}

\begin{solution}{exo:b2:cell-differentiation:8}
Con $k = 0.6$: para $m = 1$, la degradación $0.6$ supera a la síntesis
$0.52$; para $m = 0.5$, $0.30 > 0.22$; para $m = 2$, $1.2 > 0.82$: la
recta queda por encima de la curva en todo el intervalo más allá del
estado bajo, que es el único estado estacionario. Una célula cuyo \omterm{def:b2:cell-differentiation:myogenesis}{MyoD} se
degrada demasiado deprisa no puede mantener el estado encendido: recae en
el estado indiferenciado sea cual sea la inducción --- no hay músculo.
\end{solution}

\begin{solution}{exo:b2:cell-differentiation:9}
(a) Sin \omterm{def:b2:cell-differentiation:myogenesis}{mioblastos}, sin músculo esquelético: los dos factores son
redundantes para la \omterm{def:b2:cell-differentiation:differentiation}{determinación}, de modo que deben faltar ambos.
(b) Los \omterm{def:b2:cell-differentiation:myogenesis}{mioblastos} se forman y se alinean, pero nunca se fusionan ni
fabrican proteínas musculares: la miogenina es necesaria para la
\omterm{def:b2:cell-differentiation:differentiation}{diferenciación}. (c) Casi normal, porque Myf5 sustituye a \omterm{def:b2:cell-differentiation:myogenesis}{MyoD} en la
\omterm{def:b2:cell-differentiation:differentiation}{determinación}.
\end{solution}

\begin{solution}{exo:b2:cell-differentiation:10}
En un fibroblasto, los genes musculares están en una cromatina que \omterm{def:b2:cell-differentiation:myogenesis}{MyoD}
puede abrir; en una célula hepática pueden estar empaquetados en
heterocromatina, o los propios \omterm{prop:b2:cell-differentiation:transcriptional}{reguladores maestros} del hígado pueden
reprimirlos, de modo que la célula no es competente. Comprobación: tratar
células hepáticas con un fármaco que relaje la cromatina (un inhibidor de
las histona desacetilasas) antes de añadir \omterm{def:b2:cell-differentiation:myogenesis}{MyoD}, o añadir \omterm{def:b2:cell-differentiation:myogenesis}{MyoD} junto con
un factor que desplace el programa hepático, y medir la expresión de
miosina.
\end{solution}

\begin{solution}{exo:b2:cell-differentiation:11}
Las descargas sostenidas de baja frecuencia elevan el calcio intracelular
durante periodos prolongados; las vías activadas por el calcio encienden
los genes de la miosina lenta, de la biogénesis mitocondrial, de la
mioglobina y del crecimiento capilar, y apagan las isoformas rápidas. El
programa maestro de la fibra (el estado de \omterm{def:b2:cell-differentiation:myogenesis}{MyoD}, la arquitectura del
\omterm{def:b2:cell-differentiation:fibre}{sarcómero}) no se toca: el entrenamiento cambia cuáles de los genes
musculares se expresan, no la identidad de la célula, y ninguna señal de
un nervio puede reabrir el programa del cartílago, que se cerró en la
\omterm{def:b2:cell-differentiation:differentiation}{determinación}.
\end{solution}

\begin{solution}{exo:b2:cell-differentiation:12}
El estado encendido de un factor autoactivador es un punto estable de un
bucle de retroalimentación, mantenido por una síntesis continua, no por
ningún cambio en los genes; las marcas de la cromatina añaden una segunda
capa que mantiene cerrados los genes silenciados. Reiniciar exige, por
tanto, llevar el circuito por debajo de su umbral y reabrir la cromatina
a la vez, lo que el citoplasma del huevo o cuatro factores solo logran en
una pequeña fracción de las células: es posible, porque no se ha perdido
nada, pero poco eficaz, porque hay que forzar a la vez dos cerraduras
independientes.
\end{solution}

\begin{solution}{pb:b2:cell-differentiation:1}
\textbf{1.} $0.2/2.5\times 10^{-6} = \num{80000}$ \omterm{def:b2:cell-differentiation:fibre}{sarcómeros};
$0.2/25\times 10^{-6} = \num{8000}$ núcleos.
\textbf{2.} Unos \num{8000} por fibra; $4\times 10^{9}$ para el
músculo.
\textbf{3.} $\pi(30\times 10^{-6})^{2}\times 0.2 = \qty{5.7e-10}{m^3}$
por fibra; $2.8\times 10^{-4}$ \unit{m^3}, \qty{0.28}{L}, para el
músculo.
\textbf{4.} $5.7\times 10^{-10}/8000 = \qty{7e-14}{m^3} =
\qty{70000}{\micro m^3}$: el equivalente de 35 células ordinarias por
núcleo.
\textbf{5.} El interior está lleno de \omterm{def:b2:cell-differentiation:fibre}{miofibrillas} que deben recorrerlo sin
interrupción de un extremo a otro; los núcleos en el borde dejan continua
la red y quedan cerca de la membrana y de sus señales.
\textbf{6.} $4\times 10^{9}/2000 = 2\times 10^{6}$: 21 duplicaciones, 10,5
días.
\textbf{7.} $0.4m = m^{2}/(1 + m^{2}) + 0.02$. Con $m = 0.055$: izquierdo
$0.022$, derecho $0.003 + 0.02 = 0.023$.
\textbf{8.} Con $m = 2.1$: izquierdo $0.84$, derecho $0.815 + 0.02 = 0.835$.
\textbf{9.} Con $m = 0.41$: izquierdo $0.164$, derecho $0.144 + 0.02 = 0.164$.
Justo por debajo, la síntesis es menor que la degradación, de modo que $m$
baja hacia $0.055$; justo por encima, la síntesis supera a la degradación y
$m$ sube hacia $2.1$: todo desplazamiento crece --- inestable.
\textbf{10.} $0.6 > 0.41$: la primera célula pasa al estado encendido y se
diferencia; $0.3 < 0.41$: la segunda vuelve a $0.055$ y no se diferencia.
\textbf{11.} Con $b = 0.4$ la curva de síntesis parte de $0.4$ y se mantiene
por encima de la recta $0.4m$ hasta $m \approx 3.3$: el corte bajo ha
desaparecido. La célula se diferencia espontáneamente, sin inducción.
\textbf{12.} La división reduce $m$ a la mitad; las hijas heredan $2.1/2 =
1.05$, por encima del umbral $0.41$, y vuelven a subir hasta $2.1$: el
estado se restablece en cada hija. El compromiso sobrevive mientras el
nivel reducido a la mitad quede por encima del umbral.
\textbf{13.} $4$ por mm $\times$ \qty{200}{mm} $= 800$ células satélite por
fibra; hay que sustituir el \qty{10}{\%} de \num{8000} núcleos, es decir,
\num{800}.
\textbf{14.} El segmento lesionado contiene $80$ células satélite; $800/80 =
10$, de modo que 4 duplicaciones ($\times 16$, lo que da \num{1280}: 800 se
fusionan y 480 quedan para restablecer la reserva) --- \qty{72}{h}, tres
días.
\textbf{15.} Tres días de división frente a las dos o tres semanas
observadas: el resto corresponde a la eliminación de los restos por los
macrófagos, el retraso de activación, la fusión, el ensamblaje de las
\omterm{def:b2:cell-differentiation:fibre}{miofibrillas}, la reinervación y la maduración del nuevo segmento.
\textbf{16.} \omterm{def:b2:cell-differentiation:myogenesis}{MyoD} (y la miogenina) deben volver a encender los genes
musculares en las células activadas; la división y la \omterm{def:b2:cell-differentiation:differentiation}{diferenciación} son
excluyentes --- Id debe disminuir y el ciclo detenerse antes de que la
miogenina pueda actuar y se produzca la fusión.
\textbf{17.} $0.95^{n} = 0.5$: $n = 13.5$, unas catorce lesiones.
\textbf{18.} $800\times 0.95^{20} = 800\times 0.36 \approx 290$.
\textbf{19.} Cada contracción desgarra fibras, cada desgarro recurre a la
reserva, y la reserva disminuye unos pocos puntos porcentuales cada vez;
al cabo de años se agota, las fibras desgarradas ya no se reconstruyen y
los fibroblastos rellenan los huecos con colágeno.
\textbf{20.} $(80/60)^{2} = 1.78$; núcleos $8000\times 1.78 = \num{14200}$:
unos \num{6200} adicionales aportados por las células satélite.
\textbf{21.} $0.28\times 1.78 = \qty{0.5}{L}$.
\textbf{22.} Ha cambiado: los genes de las isoformas de la miosina, la
densidad mitocondrial y capilar, la mioglobina. No ha cambiado: la
identidad muscular (el estado de \omterm{def:b2:cell-differentiation:myogenesis}{MyoD}), la arquitectura del \omterm{def:b2:cell-differentiation:fibre}{sarcómero},
los núcleos y su genoma.
\textbf{23.} Las fibras son posmitóticas y no pueden multiplicarse; las
únicas vías son fibras más gruesas y más núcleos por fibra aportados por
las células satélite.
\textbf{24.} Entrenamiento: un engrosamiento limitado, puesto que la fibra
no puede añadir núcleos; lesión: ninguna reparación --- los segmentos
destruidos son sustituidos por tejido fibroso, como en la distrofia.
\textbf{25.} \num{8000} núcleos por fibra; umbral $m \approx 0.41$;
unos tres días de división de las células satélite para una reparación.
\end{solution}
